% Lösungen Einheit 1 von 5 – Prozente als Anteile (zusammenbau v0.1)
\einheitenkopf{Einheit 1 von 5 · Prozente als Anteile}

%% TODO Lösung der Erklärzeile 9f) fehlt in der Bank
\erg{9}{a) $50\,\%$ \quad b) $25\,\%$ \quad c) $60\,\%$ \quad d) $30\,\%$ \quad e) $90\,\%$ \quad f) \ldots \quad g) $37\,\%$ \quad h) $\frac{55}{100} = 55\,\%$ \quad i) $64\,\%$ \quad j) $\frac{1}{100} = 1\,\%$ \quad k) $100\,\% - 35\,\% = 65\,\%$ \quad l) $5$ von $100$ Paketen kommen zu spät. ($5\,\% = \frac{5}{100}$; „jedes 5.“ wären $20\,\%$) \quad m) Aussage 2 ist falsch: $200 : 600 \approx 0{,}33$, also ein Drittel. Richtig: Ein Drittel der Kinder fährt mit dem Rad. (Aussage 1 stimmt: $150 : 600 = 0{,}25$.) \quad n) Biogas: $100 - 38 - 24 - 9 - 3 = 26$, also $26\,\%$}
\erg{10}{a) $\frac{1}{4} < \frac{3}{10} < 35\,\% < 0{,}4$ (in Prozent: $\frac{1}{4} = 25\,\%$, $\frac{3}{10} = 30\,\%$, $35\,\%$, $0{,}4 = 40\,\%$)}
\erg{11}{a) Fehler: Tom liest den Nenner $25$ als Prozentzahl. Richtig: $\frac{1}{25} = \frac{4}{100} = 4\,\%$. \quad b) Mit $25$ erweitert ist $\frac{1}{4} = \frac{25}{100}$, und Prozent heißt Hundertstel. \quad c) $3$ der $10$ Teile gefärbt \quad d) 7a: $\frac{18}{25} = 72\,\%$; 7b: $\frac{21}{30} = \frac{7}{10} = 70\,\%$ – in der 7a ist der Anteil größer.}
